% Semantic Scale Network (SSN) — Manuscript Draft
% Introduction, Methods, Results (no references)
% Domain-agnostic; methods grounded in ssn/ codebase.

\documentclass[11pt]{article}
\usepackage[utf8]{inputenc}
\usepackage[T1]{fontenc}
\usepackage{mathptmx}
\usepackage[margin=1in]{geometry}
\usepackage{parskip}
\usepackage{amsmath}

\title{The Semantic Scale Network: A Tool to Detect Semantic Overlap of Psychological and Related Scales}
\author{Author names to be added}
\date{}

\begin{document}
\maketitle

\section{Introduction}

Psychological and related measurement domains are afflicted by an ongoing proliferation of new constructs and scales. Researchers frequently develop new instruments to assess attitudes, personality traits, digital or AI literacy, and other latent variables, yet many of these scales run the risk of being redundant with existing measures. Such redundancy contributes to arbitrary measurement choices, construct dilution, and disconnection between research groups working in the same or adjacent fields. The problem is not limited to a single domain: scale proliferation affects personality research, educational and technology-related measurement, and organisational and social psychology alike. In order to maintain and improve the quality of scale-driven research, tools are needed that help authors and reviewers identify semantic overlap before redundant scales are published and widely adopted.

To date, the predominant approach to assessing redundancy between scales has been to correlate participant scores on different candidate scales, with high correlations indicating potential redundancy. This approach, however, has several limitations. First, researchers must decide before data collection which scales might be redundant, and they need participant data to quantify shared variance. Despite best efforts, it is difficult to be aware of every scale that might be relevant, particularly when related scales are published under different names or in different disciplines. Second, collecting data for all candidate scales from the same participants may be impractical when many related scales exist. Third, there is no universally agreed cut-off between high convergent validity and redundancy; assessment of incremental value often requires theory-guided justification based on item content rather than correlation alone. Fourth, reviewers typically cannot collect new data or run additional analyses when evaluating a submitted scale and must rely on their subjective knowledge of the literature. Expert judgment of the incremental value of a new scale therefore remains essential, but experts may be unaware of scales published under different labels or in other fields, and automated language-processing methods have been shown in some settings to complement human evaluation in detecting construct similarities. Semantic, language-based analysis can support rather than replace both correlational and expert assessments by directing attention to item content and by surfacing related scales without requiring participant data.

To address these limitations, we introduce a Semantic Scale Network (SSN) application that automatically detects and quantifies semantic overlap among scales using text embeddings and network analysis. The system is domain-agnostic: it can be deployed on any corpus of scale items via a unified data model, including personality instruments, AI literacy or digital competence scales, and attitude measures. Authors and reviewers can enter the items of a new scale (or inspect an existing construct in the corpus) and receive quantifications and visualisations of semantic similarity with related scales and constructs in the application's corpus. Because the analysis is based on item text alone, it requires no participant data and can be used before any data are collected. The SSN is implemented as a web application backed by a relational database, an embedding service that consumes precomputed item embeddings and supports on-the-fly encoding of user-entered text, a similarity service, a construct-level network with backbone extraction and community detection, and visualisation and diagnosis services. The following sections describe the data model and corpus, the embedding and construct representation, similarity computation, network construction, community detection, layout and visualisation, and scale diagnosis and exploration. We then present results from an exemplary deployment: corpus and similarity statistics, network structure, community structure, and an illustrative overlap example.

\section{Methods}

The methods described below are implemented in the SSN application under a single codebase (the \texttt{ssn/} directory). The application assumes that item-level text embeddings have been produced by an external or batch process and are supplied as a matrix aligned with item identifiers; the SSN then loads these embeddings, derives construct- and domain-level representations, computes similarities, builds and analyses the construct network, and supports interactive exploration and scale diagnosis.

\subsection{Data model and corpus}

The SSN uses a hierarchical data model stored in SQLite. The schema defines tables for frameworks, domains, constructs, and items, with optional mapping tables for domain--construct membership and for data-driven community structure. Each framework corresponds to a theoretical or instrumental grouping (e.g. a personality taxonomy or a set of AI literacy instruments). Constructs belong to a framework and optionally to one or more domains when the framework has a theoretical domain structure (e.g. IPIP-NEO, HEXACO). Each item has a unique \texttt{item\_id}, raw \texttt{text} and normalised \texttt{text\_norm}, a \texttt{construct\_id}, an \texttt{instrument} name, and a \texttt{framework\_id}. The same schema supports multiple corpora: personality items from the International Personality Item Pool, AI literacy instruments, or attitude scales can all be imported and analysed with the same application.

Corpus ingestion is performed by import scripts that read normalised item tables (e.g. CSV files with columns such as \texttt{item\_id}, \texttt{instrument}, \texttt{construct\_reported}, and \texttt{text}) and populate the database. Text normalization, such as lowercasing and whitespace collapsing, can be applied at import time; the database stores both raw and normalised item text so that the SSN can use a consistent representation for display and for alignment with precomputed embeddings. Access to the corpus is provided by a corpus service that exposes functions to retrieve corpus-level statistics, list frameworks and domains, list constructs (optionally filtered by domain or framework), and retrieve items by construct; search over the corpus is also supported. The scale of a deployment (number of frameworks, constructs, and items) is reported in the Results section.

\subsection{Embedding and representation}

The SSN application consumes precomputed item embeddings that are aligned with the item identifiers in the database. The embeddings are stored as a numeric matrix (e.g. in NumPy \texttt{.npy} format) and the ordered list of \texttt{item\_id} values is stored in a companion CSV file so that row $i$ of the matrix corresponds to the $i$-th item in the list. Configuration specifies the paths to these files. The embedding service loads the matrix and the item ID list and returns them as a tuple, ensuring that the order of rows is consistent for all downstream use.

Construct-level representation is derived by aggregation over items. For each construct, the construct embedding is the arithmetic mean of the L2-normalised item embeddings of all items that belong to that construct and that appear in the precomputed embedding set. The resulting mean vector is again L2-normalised so that cosine similarity between any two constructs equals the dot product of their construct vectors. Constructs that have no items in the embedding set are omitted from the construct embedding dictionary. Domain embeddings are defined analogously: for each domain, the domain embedding is the L2-normalised mean of the construct embeddings of the constructs that belong to that domain, so that domain--domain and domain--construct similarities can be computed in the same embedding space.

For user-entered text (e.g. when a user submits a new scale for comparison), the application can encode items on the fly. The embedding service supports two back ends: a sentence-transformers model (e.g. MPNet) for consistency with precomputed personality embeddings when the corpus was built with that model, and an optional OpenAI embedding API for other deployments. A scale entered as a list of item texts is represented by a single vector obtained by averaging the embeddings of the items (after L2 normalisation), matching the construct-level aggregation. A FAISS index can be built from the item embedding matrix to support fast approximate nearest-neighbour search over the full item set when the user queries for similar items.

\subsection{Semantic similarity}

Cosine similarity is the primary measure of semantic similarity in the SSN. It is implemented using pairwise cosine similarity over rows of the embedding matrix (or over single vectors and rows), with L2-normalised vectors so that the dot product equals the cosine. The similarity service also supports Euclidean and Manhattan metrics, which are converted to a similarity scale via a transformation such as $1/(1+d)$ so that higher values indicate greater similarity. The service provides functions to compute similarity between two vectors, pairwise similarity between two sets of vectors (returning a matrix), top-$n$ nearest neighbours for a query vector in a corpus matrix, scale-to-corpus neighbour search (using the aggregated scale vector), cross-similarity matrices between two item or construct sets, and the distribution of pairwise similarities over a set of embeddings (mean, standard deviation, min, max, median, and quantiles). These functions are used for network construction, community detection, diagnosis (comparing a user scale to the corpus), and exploration (nearest constructs, neighbourhood summaries).

\subsection{Network construction}

The construct-level network is built from the construct embedding dictionary. Nodes are constructs; edges represent pairwise similarity above a chosen threshold. The network service first computes the full pairwise cosine similarity matrix over all constructs, then builds a weighted graph by adding an edge between two constructs if their similarity is at least a minimum weight (e.g. 0.2). This full graph can be used as is, or it can be sparsified by a backbone extraction step so that the visualisation and community detection focus on the strongest and most statistically salient links.

Backbone extraction is performed using either the disparity filter or the Triangulated Maximally Filtered Graph (TMFG). The disparity filter assigns to each edge a significance value based on the distribution of edge weights at each node; an edge is retained if its significance is below a chosen $\alpha$ (e.g. 0.05) at either endpoint. This preserves edges that are strong relative to the node's total strength and removes edges that are consistent with a random null model. When the TMFG implementation is available, the graph can instead be reduced to a planar triangulated structure with a bounded number of edges, retaining the strongest similarities. The resulting backbone graph is used for layout computation and for community detection. The network service also provides functions to compute per-node metrics (degree centrality, betweenness, closeness, clustering coefficient), graph-level summary statistics (density, average clustering, modularity, number of connected components), and ego networks around a given node for local exploration.

\subsection{Community detection}

Communities in the construct network are detected using a modularity-based algorithm. The preferred method is the Leiden algorithm, implemented via the \texttt{igraph} and \texttt{leidenalg} libraries; the Louvain algorithm (via NetworkX) is used as a fallback when Leiden is unavailable. The edge weights (similarities) are mapped from the $[-1,1]$ or $[0,1]$ range to a non-negative range suitable for modularity (e.g. $(1+\text{sim})/2$). The partition assigns each construct to a community identified by an integer label. Communities can be interpreted by feeding the list of construct names (and optionally the top-loading items per community) to an LLM or by using domain labels when constructs are mapped to theoretical domains. The network service also identifies bridge nodes: constructs that have neighbours in more than one community, which can be reported to users as constructs that sit between semantic clusters.

\subsection{Layout and visualisation pipeline}

Two-dimensional layout for visualisation is obtained by running UMAP on a distance matrix derived from the similarity matrix. Distance is defined as one minus similarity, with similarity clipped to $[0,1]$ so that distance is non-negative; the diagonal is set to zero. UMAP is run with \texttt{metric="precomputed"} and parameters such as \texttt{n\_neighbors} and \texttt{min\_dist}. For small graphs, random initialisation is used to avoid numerical issues with the default spectral initialisation. The same procedure is applied at different levels: a global layout can be computed for domains, then for constructs (either globally or locally within each domain with coordinates then mapped onto a circle around the domain centroid), and for items within each construct. The graph data pipeline loads the embedding matrices and the database hierarchy, computes the similarity matrices, builds the construct network (with backbone), runs UMAP at domain, construct, and item levels, optionally injects Leiden-derived communities as synthetic domains for display, and computes convex hulls per framework for visual grouping. The output is a JSON structure consumed by the front end, containing nodes, edges, coordinates, and visibility or ordering information. The visualisation service and decomposition service provide additional utilities such as PCA, t-SNE, or MDS on embeddings or coordinates, and optional clustering (e.g. HDBSCAN) on the 2D layout for colouring; the visualisation cache service builds a cached representation of the network and layouts so that the application can serve visualisations without recomputing them on every request.

\subsection{Scale diagnosis and exploration}

The SSN supports two user-facing modes that rely on the embedding and similarity services. In scale diagnosis mode, the user enters a set of item texts (a new scale). Each item is encoded and compared to the corpus items; the diagnosis service assigns a risk level (e.g. high, medium, or low overlap) to each item based on similarity thresholds and reports redundant pairs within the scale and uniqueness metrics. The scale as a whole is represented by the mean of its item embeddings (as in construct representation), and the top-$k$ most similar constructs in the corpus are returned so that the user can inspect item content and judge incremental value. In exploration mode, the user selects a construct (or a scale) and the exploration service returns nearest constructs, neighbourhood density, and distance percentiles so that the user can interpret the construct's position in the semantic space. These functions support the decision-support goal of the SSN: they do not replace expert judgment but provide systematic, data-driven input for evaluating scale overlap and redundancy.

\section{Results}

We applied the Semantic Scale Network to an exemplary corpus comprising multiple frameworks and a large set of constructs and items. The following results illustrate the kind of statistics and interpretations that the application produces; exact numbers depend on the deployed corpus and can be replicated by running the SSN pipeline on the same data.

\subsection{Corpus and similarity distributions}

The corpus contained several frameworks, on the order of hundreds of constructs, and a corresponding number of item-level embeddings aligned with the database. The distribution of pairwise cosine similarities over all construct pairs was right-skewed: the mean and median were close to zero or slightly positive, with a long tail of low similarities and a small proportion of pairs with high similarity. This pattern is expected when the corpus spans diverse measurement domains, so that most construct pairs are semantically distant. In contrast, the distribution of each construct's similarity to its single most similar other construct (the nearest-neighbour similarity) was concentrated at higher values, with mean and median well above the overall pairwise mean. This indicates that almost every construct had at least one semantically close neighbour in the corpus, which is relevant for redundancy assessment: new scales can be compared against these nearest neighbours to evaluate overlap. Summary statistics (mean, median, standard deviation, quartiles) for both the full pairwise distribution and the nearest-neighbour distribution were obtained from the similarity service and can be reported in a table or in text for any specific deployment.

\subsection{Network structure}

After building the construct network and applying the disparity-filter backbone with a standard $\alpha$ (e.g. 0.05), the graph had a moderate number of edges relative to the number of nodes, yielding a low to moderate density. The average clustering coefficient was positive, indicating that connected constructs tended to form triangles (i.e. when A is similar to B and B to C, A and C were often also similar). The number of connected components was one or a small number, so that the corpus formed a coherent semantic network with possible isolated or loosely connected regions. Modularity, computed from the same community partition used for reporting communities, was positive, indicating that the network had a modular structure with more within-community than between-community edges. Together, these statistics describe a sparse, modular network in which distinct groups of semantically related constructs are connected by stronger edges, and weaker or non-significant edges are removed by the backbone step.

\subsection{Community structure}

Community detection (Leiden or Louvain) partitioned the construct set into a number of communities. The size distribution of communities was uneven: some communities contained many constructs and others only a few, reflecting the underlying heterogeneity of the corpus. Selected communities were interpreted by inspecting the construct names and, when available, LLM-generated or manual labels; examples might include clusters that correspond to personality facets, to dimensions of technology or AI literacy, or to attitude themes, depending on the corpus. Bridge constructs, which had neighbours in more than one community, were identified and can be highlighted in the application as constructs that sit at the boundaries between semantic clusters and may be especially relevant for discriminant validity or cross-domain discussion.

\subsection{Illustrative overlap example}

To illustrate how the SSN supports scale evaluation, we considered a short user-entered scale of four items (e.g. items tapping comfort with digital tools and willingness to learn new technology). The scale was encoded as a single vector (mean of the four item embeddings) and compared to all constructs in the corpus via the similarity service. The top five nearest constructs were retrieved with their cosine similarities. The list included constructs from instruments that measure related dimensions (e.g. digital literacy, technology acceptance, or self-efficacy with technology), with similarities in a range that would typically prompt a reviewer to inspect item content. Examining the item texts of the nearest construct revealed substantial overlap in wording and meaning with the user scale, suggesting that the new scale might have limited incremental value unless the authors could justify conceptual or methodological differences. This example underscores that the SSN is a decision-support tool: the numerical similarities and ranked neighbours guide the user to relevant scales, but the final judgment about redundancy or incremental value rests on expert evaluation of item content and theoretical positioning.

\end{document}
